\section{The exchange criterion and criticality}\label{sec:exchange}
The same wedge identities describe the sum-of-squares problem exactly
at each fixed order, and show what the margins of positive certificates
can be.

Fix $n\ge3$ and use one-based depths $1\le p<n$:
\[
 w_p^-=w_{p-n},\qquad w_p^+=w_{n-p},\qquad
 u_{pq}=w_p^-\wedge w_q^-,\quad v_{pq}=w_p^+\wedge w_q^+,\quad
 z^{\rm mix}_{pq}=w_p^-\wedge w_q^+.
\]
Both pure sides are oriented by increasing depth, so $v_{pq}$ is the
negative of the wedge ordered by increasing label. Unlike the stabilized
corner, these variables include all noncentral diagonals.

The central symbol splits off as explicit positive squares:
\begin{equation}\label{eq:center}
 F_n=F_n|_{w_0=0}+
       \sum_{\ell\ne0}2n(n-|\ell|)z_{0\ell}^2.
\end{equation}
Indeed the central shift is the identity, commutes with every shift,
and is Frobenius orthogonal to the other shifts. Therefore $F_n$ is
SOS if and only if its noncentral restriction is SOS.

\subsection{The pure budget}
On $W_n=\Lambda^2\R^{n-1}$ define
\begin{equation}\label{eq:budget-baseline}
 \begin{split}
 D_n&=\diag(2pq)_{1\le p<q<n},\\
 (J_n)_{pq,rs}&=\one_{q-p=s-r}\min(p,r,n-q,n-s),\qquad p<q,\ r<s.
 \end{split}
\end{equation}
Let $\calP_{n-1}$ be the pure Pl\"ucker span. On mixed coordinates put
\begin{equation}\label{eq:budget-mixed}
 (M_n^0)_{pq,rs}=2pq\one_{(p,q)=(r,s)}
 -2\one_{p+q=r+s}\min(p,q,r,s,n-p,n-q,n-r,n-s).
\end{equation}
For a symmetric matrix $\Xi$ on $W_n$, the realignment $\calR(\Xi)$ is
the mixed matrix \eqref{eq:realignment-entries}. Let
\begin{equation}\label{eq:exchange-set}
 \begin{split}
 \mathcal L_n&=\{\Xi=\Xi^T:\calR(\Xi)s_g=0, 2\le g\le n\},\\
 \mathcal C_n&=\{\Xi\in\mathcal L_n:M_n^0+\calR(\Xi)\succeq0\},
 \qquad s_g=\sum_{p+q=g}\mathbf e_{pq}.
 \end{split}
\end{equation}
Define the pure-completion norm by
\begin{equation}\label{eq:pure-norm}
 \nu_n(Y)=\max_{\substack{Z\succeq0, Z\perp\calP_{n-1}\\
                         \tr(D_nZ)=1}}
                  \tr|Z^{1/2}YZ^{1/2}|,
\end{equation}
where $|Y|$ is the spectral absolute value, and define the budget
\begin{equation}\label{eq:gamma}
                       \gamma_n=\min_{\Xi\in\mathcal C_n}\nu_n(2J_n+\Xi).
\end{equation}
The normalizations in this section differ from the zero-based corner
only by the stated index shift and by naming the positive minimum
kernel $J_n$, so that the pure cross block is $-(2J_n+\Xi)$.

\begin{theorem}[Exact exchange criterion]\label{thm:exchange}
For every $n\ge3$, the maximum in \eqref{eq:pure-norm} and the minimum in
\eqref{eq:gamma} are attained, $\nu_n$ is a norm, and
\[
                         F_n\text{ is SOS}\quad\Longleftrightarrow\quad\gamma_n\le1.
\]
Equivalently, the latter condition is the existence of
$\Xi\in\mathcal C_n$ and $E\in\calP_{n-1}$ such that
\begin{equation}\label{eq:budget-feasible}
                    D_n+E\pm(2J_n+\Xi)\succeq0.
\end{equation}
If $\gamma_n<1$, such a certificate can be chosen rational, with both
pure blocks positive definite and the mixed block positive definite on
the complement of the forced means.
\end{theorem}
\begin{proof}
The complete wedge-frame and symmetry argument of
Lemma~\ref{lem:corner} applies to the noncentral restriction.
For the signed shifts $S_{p-n}$ and $S_{n-q}$, multiplication and
cancellation of the overlapping strings give
\[
 \begin{split}
 &\langle[S_{p-n},S_{n-q}],[S_{r-n},S_{n-s}]\rangle_F\\
 &\qquad=2\one_{p-q=r-s}
       \min(p,q,r,s,n-p,n-q,n-r,n-s).
 \end{split}
\]
Indeed different offsets are orthogonal; on a common offset the two
remaining end strings have this common length. The norm terms have
weights $2pq$. Applying \eqref{eq:plucker-exchange} to each raw mixed
pair with equal offset sends it to a mixed pair with equal total and
a pure pair with equal gap. This gives $M_n^0$ and $J_n$ above, and the
averaged Gram is
\[
 \begin{pmatrix}
 D_n+E&-(2J_n+\Xi)&0\\
 -(2J_n+\Xi)&D_n+E&0\\
 0&0&M_n^0+\calR(\Xi)
 \end{pmatrix}.
\]
The cross sign follows from
$u_{pq}v_{rs}-z^{\rm mix}_{pr}z^{\rm mix}_{qs}+z^{\rm mix}_{ps}z^{\rm mix}_{qr}=0$: specifying the parallel
mixed entry to be $\Xi_{pq,rs}$ forces the pure cross entry
$-2(J_n)_{pq,rs}-\Xi_{pq,rs}$. All pure four-forms remain free.
The sum and difference of the pure sides give
\eqref{eq:budget-feasible}.

For totals $g\le n$, the row-sum identity
\[
 \sum_{r=1}^{g-1}\min(p,g-p,r,g-r)=p(g-p)
\]
makes $M_n^0$ a weighted complete-graph Laplacian. For $g>n$, reversal
$p\mapsto n-p$, $q\mapsto n-q$ gives the block at $2n-g$ plus
$2n(g-n)I$. Hence $M_n^0\succeq0$ and its kernel is exactly the span
of $s_2,\ldots,s_n$. The least-coefficient argument of
Lemma~\ref{lem:means} forces these same means for every positive Gram.
Thus imposing $\mathcal L_n$ loses no feasible matrix.

The set $\mathcal C_n$ contains $0$ and is closed and convex.
Its mixed diagonal is fixed. Positive semidefiniteness bounds every
entry by the geometric mean of its diagonal entries, and the parallel
entries recover every entry of $\Xi$. This proves compactness.

For a fixed symmetric $Y$, consider
\begin{equation}\label{eq:pure-primal}
 \inf\{\tau:\exists E\in\calP_{n-1},\quad \tau D_n+E\pm Y\succeq0\}.
\end{equation}
It is finite and nonnegative. In a bounded sublevel set, the diagonals
of both positive matrices are bounded, since $E$ has zero diagonal;
positive semidefiniteness bounds their off-diagonal entries and hence
$E$. Closedness gives attainment. Taking $E=0$ and $\tau$ sufficiently
large gives strict feasibility. Finite-dimensional semidefinite duality
identifies \eqref{eq:pure-primal} with
\[
 \max\{\langle Y,Z_+-Z_-\rangle:
 Z_\pm\succeq0,\ Z_++Z_-\perp\calP_{n-1},\
 \langle D_n,Z_++Z_-\rangle=1\}.
\]
This dual set is nonempty and compact. For fixed $Z=Z_++Z_-$, the
maximum is $\tr|Z^{1/2}YZ^{1/2}|$: on the range of $Z$, the difference
$Z_+-Z_-$ is $Z^{1/2}UZ^{1/2}$ with $-I\preceq U\preceq I$.
The argument includes singular $Z$. It proves \eqref{eq:pure-norm} and its
attainment. The formula is homogeneous and subadditive; it separates
nonzero $Y$ because a positive multiple of the identity is an admissible
$Z$. Thus it is a norm.

Continuity and compactness give attainment in \eqref{eq:gamma}.
An attained value $\tau\le1$ in \eqref{eq:pure-primal} can be increased
to $1$ by adding $(1-\tau)D_n\succeq0$. This proves the equivalence,
including equality at $\gamma_n=1$.

Finally, if $\gamma_n<1$, move $\Xi$ a sufficiently small distance
toward $0$. The mixed block becomes positive definite on the forced
kernel complement and the pure budget remains below $1$. Rational
points are dense in the rational subspaces $\mathcal L_n$ and
$\calP_{n-1}$. Approximate within the resulting strict inequalities.
This gives the rational certificate claimed.
\end{proof}

\begin{theorem}[Criticality of the pure budget]\label{thm:criticality}
For every $n\ge3$, $\tfrac12\le\gamma_n\le4$. For every $n\ge128$,
\begin{equation}\label{eq:criticality}
                  \gamma_n\ge1-\frac{33\lceil\log_2n\rceil^2}{n^2}.
\end{equation}
If an SOS certificate has pure blocks
$D_n+E\pm(2J_n+\Xi)\succeq\mu D_n$, then
$\mu\le1-\gamma_n$. In particular its relative pure margin is at most
the displayed deficit when $n\ge128$.
\end{theorem}
\begin{proof}
Appendix~\ref{app:criticality} proves the quantitative bounds from an
explicit decomposable wedge test and the complete low means.
For the margin assertion, subtract $\mu D_n$ from both pure blocks.
The same $E,\Xi$ is then feasible at budget $1-\mu$, so
$\gamma_n\le1-\mu$.
\end{proof}

\subsection{Consequences for certificates}
Any feasible certificates at increasing orders must therefore lose their
relative pure margin at least at the rate imposed by
\eqref{eq:criticality}. This is the precise meaning of criticality here.
The theorem gives a lower bound for $\gamma_n$, not convergence to $1$.
Indeed Theorem~\ref{thm:negative} and the exact criterion imply
$\gamma_n>1$ throughout the large-order non-SOS range.
